\documentclass[10pt]{article}
\usepackage[margin=1.8cm,top=1.5cm,bottom=1.8cm]{geometry}
\usepackage[utf8]{inputenc}
\usepackage[T1]{fontenc}
\usepackage{lmodern}
\usepackage{booktabs}
\usepackage{graphicx}
\usepackage{xcolor}
\usepackage{framed}
\usepackage{amsmath}
\usepackage{multicol}
\usepackage{hyperref}
\usepackage{array}
\usepackage{colortbl}
\usepackage{calc}

\definecolor{khpblue}{RGB}{0,91,148}
\definecolor{riskred}{RGB}{180,0,0}
\definecolor{safegreen}{RGB}{0,130,0}
\definecolor{warnamber}{RGB}{200,100,0}
\definecolor{lightgray}{RGB}{248,248,248}
\definecolor{darkgray}{RGB}{80,80,80}
\definecolor{shadecolor}{RGB}{240,245,252}

\hypersetup{colorlinks=true,linkcolor=khpblue,citecolor=khpblue,urlcolor=khpblue,
  pdftitle={Cassandre Grand Composite -- KHP Hackathon 2026},pdfauthor={Team 055}}

% Section formatting without titlesec
\makeatletter
\renewcommand{\section}{\@startsection{section}{1}{\z@}%
  {-6pt}{3pt}{\large\bfseries\color{khpblue}\hrule\vspace{2pt}}}
\renewcommand{\subsection}{\@startsection{subsection}{2}{\z@}%
  {-4pt}{2pt}{\normalsize\bfseries\color{darkgray}}}
\renewcommand{\subsubsection}{\@startsection{subsubsection}{3}{\z@}%
  {-3pt}{1pt}{\small\bfseries\itshape}}
\makeatother

\setlength{\parskip}{3pt}
\setlength{\parindent}{0pt}
\setlength{\columnsep}{1.5em}

\pagestyle{myheadings}
\markboth{Team 055 --- KHP Guardrail Submission, March 2026}{}

% Coloured framed box using xcolor + framed
\newenvironment{bluebox}{%
  \def\FrameCommand{{\color{khpblue}\vrule width 3pt}\colorbox{shadecolor}}%
  \MakeFramed{\advance\hsize-\width\FrameRestore}}%
{\endMakeFramed}

\newenvironment{redbox}{%
  \def\FrameCommand{{\color{riskred}\vrule width 3pt}\colorbox{red!8}}%
  \MakeFramed{\advance\hsize-\width\FrameRestore}}%
{\endMakeFramed}

\newenvironment{amberbox}{%
  \def\FrameCommand{{\color{warnamber}\vrule width 3pt}\colorbox{orange!7}}%
  \MakeFramed{\advance\hsize-\width\FrameRestore}}%
{\endMakeFramed}

\begin{document}

%% ═══════════════════════════════════════════════════════════════════
%% ONE-PAGER (page 1, single column)
%% ═══════════════════════════════════════════════════════════════════
\begin{center}
{\LARGE\bfseries\color{khpblue}
Cassandre Grand Composite:\\[4pt]
A Bilingual Multi-Turn Guardrail for Youth Mental Health Risk Detection}\\[6pt]
{\large\color{darkgray} Team 055 \quad\textbar\quad KHP Hackathon 2026 \quad\textbar\quad One-Pager}
\end{center}

\vspace{4pt}
\begin{bluebox}
\small
\textbf{\color{khpblue}Problem.}\;
The Kids Help Phone (KHP) virtual assistant fails to reliably identify youth in genuine crisis.
Across 1,977 stress-test scenarios, the baseline chatbot showed critical miss rates of
97.7\% for hate-speech crises, 83.7\% for body-image crises, and 81.4\% for home violence disclosures.
For a service where a missed signal can cost a life, this is unacceptable.

\smallskip
\textbf{\color{khpblue}Method.}\;
We built \textit{Cassandre Grand Composite v3}: a hybrid guardrail combining
(1) a bilingual rule-based multi-turn analyser using a Drift Diffusion Model (DDM)
calibrated for neurodivergent youth, with
(2) a clinical-grade LLM judge guided by a C-SSRS-aligned bilingual prompt
covering English and Qu\'ebec French joual.
A context-aware scoring rule distinguishes situational frustration from genuine ideation
by detecting whether a clear external stressor is present.
Provider fallback ensures availability: OpenAI $\to$ Cohere $\to$ Mistral.

\smallskip
\textbf{\color{khpblue}Key Results (hidden-validation-102).}
\begin{center}
\begin{tabular}{lcccc}
\toprule
\textbf{Version} & \textbf{F1} & \textbf{Precision} & \textbf{Recall} & \textbf{Latency (ms/s)}\\
\midrule
Best run (v2.5) & \textbf{0.880} & \textbf{0.917} & 0.846 & 1{,}128\\
v3 (few-shot)   & 0.850 & 0.871 & 0.831 & 1{,}189\\
v3 + bilingual prompt v2 & 0.838 & 0.803 & \textbf{0.877} & 1{,}289\\
v3 + bilingual prompt v3 & --- & --- & --- & ---\\
\bottomrule
\multicolumn{5}{l}{\scriptsize v3 bilingual prompt just submitted; evaluation pending.}
\end{tabular}
\end{center}

\smallskip
\textbf{\color{khpblue}KHP Relevance.}\;
Recall of 0.877 --- the highest achieved --- means fewer missed crises.
Qu\'ebec French joual coverage addresses a critical gap:
31\% of the evaluation data is French or mixed-language.
Neurodivergent DDM profiles (ADHD/ASD) address systematic under-detection
of flat-affect and impulsive crisis disclosures.

\smallskip
\textbf{\color{riskred}Bottom line.}\;
No guardrail, however robust, justifies deploying an AI chatbot as a primary
mental health intervention for children.
This system is a harm-reduction triage layer.
The right question is not ``how good is the guardrail?'' but ``why is there not a human here?''
\end{bluebox}

\newpage

%% ═══════════════════════════════════════════════════════════════════
%% FULL REPORT (two-column)
%% ═══════════════════════════════════════════════════════════════════
\begin{center}
{\large\bfseries\color{khpblue} Full Technical Report}
\end{center}

\begin{multicols}{2}

%% ── 1. SYSTEM OVERVIEW ───────────────────────────────────────────
\section{System Overview}

\textbf{Cassandre Grand Composite v3} is a stacked hybrid input guardrail.
Given a conversation $x$ (single or multi-turn, English or French),
the composite score is:

\[
s = \begin{cases}
s_{\text{LLM}} & \text{if } s_{\text{LLM}} \geq 0.85\\[2pt]
0.4\,s_{\text{Cass}} + 0.6\,s_{\text{LLM}} & \text{otherwise}
\end{cases}
\]

\noindent The guardrail flags \texttt{HIGH\_RISK} when $s \geq 0.40$.

\subsection{Cassandre Multi-Turn Analyser}

A fully deterministic component (DDM seed = 42) scoring conversation \emph{trajectory}:
\[
s_{\text{Cass}} = 0.6\,s_{\text{DDM}} + 0.4\,s_{\text{Conv}}
\]

\textbf{ConvoAnalyzer} extracts five weighted linguistic features:
escalation trajectory (0.35), final-turn risk density (0.25),
hedge word erosion (0.15), engagement collapse (0.15),
first-to-last risk delta (0.10).
French keywords are patched in at import time via \texttt{fr\_patch.py},
extending both \texttt{\_RISK\_KEYWORDS} and \texttt{\_HEDGE\_WORDS}
with 80+ Qu\'ebec French terms
(\textit{idées noires, pu capable, chu tanné, en finir}, etc.).

\textbf{DDM} (Drift Diffusion Model) models risk as
evidence accumulation across turns \cite{ratcliff2008}:
\[
E[t{+}1] = E[t] + v\,\phi(t) + \sigma\,\varepsilon,\quad \varepsilon\sim\mathcal{N}(0,1)
\]
Three cognitive profiles adapt boundary $a$, drift $v$, noise $\sigma$, and bias $z$:
\emph{neurotypical} (default),
\emph{ADHD} (lower boundary, higher noise \cite{karalunas2012}),
\emph{ASD} (higher starting bias for flat-affect under-detection \cite{south2010}).
Profile is auto-detected from syntactic fragmentation and affect vocabulary.

\subsection{Bilingual LLM Judge}

An LLM-as-judge using a C-SSRS-aligned bilingual prompt
(14 annotated examples: English + Qu\'ebec French).
The key design innovation is a \textbf{context rule}:
identical surface language is scored differently based on whether
a clear external stressor is present
(\textit{``I give up on this essay''} $\neq$ \textit{``I give up on everything''}).
Backend: OpenAI (primary) $\to$ Cohere $\to$ Mistral.

\subsection{Pipeline}

\begin{center}\small
\texttt{Input (EN/FR)}\\
$\downarrow$\\
\texttt{FR Patch (side-effect import)}\\
$\swarrow\quad\quad\searrow$\\
\texttt{Cassandre}\quad\quad\texttt{LLM Judge}\\
$\searrow\quad\quad\swarrow$\\
\texttt{Blend + Confidence Override}\\
$\downarrow$\\
\texttt{threshold 0.40} $\to$ \texttt{HIGH/LOW RISK}
\end{center}

%% ── 2. DATA GENERATION & DEI ────────────────────────────────────
\section{Data Generation \& DEI Coverage}

\subsection{Seed Validation Set}
94 conversations: 52 low-risk (55.3\%), 42 high-risk (44.7\%).
Languages: 67\% English, 18\% French, 15\% mixed.
Categories: Social Relationships (31\%), Isolation (28\%),
Suicide (15\%), School (10\%), and six further classes.

\subsection{Synthetic Dataset}

We generated \textbf{1,977 conversations} using a DEI-stratified matrix
via Cohere (Table~\ref{tab:dei}).

\begin{center}
\captionof{table}{DEI coverage dimensions}\label{tab:dei}
\small
\begin{tabular}{ll}
\toprule
\textbf{Dimension} & \textbf{Coverage}\\
\midrule
Age & 10, 12, 14, 17, 20\\
Identity & 2SLGBTQ+, cisgender\\
Neurodivergence & Autistic/ADHD, NT\\
Ethnicity & Racialized, general\\
Register & Direct, slang, euphemistic,\\
 & fragmented\\
Language & EN, FR (QC joual), mixed\\
Risk level & Low + high, all 23 categories\\
\bottomrule
\end{tabular}
\end{center}

\textbf{Limitation:} Generated conversations carry pass/fail/critical
auto-judge labels, not binary 0/1 labels.
Used exclusively for red-team analysis, not model training.

%% ── 3. GUARDRAIL DESIGN ─────────────────────────────────────────
\section{Guardrail Design}

\subsection{Prompt Engineering Evolution}

Three prompt versions were evaluated on the hidden set:

\textbf{v1 (C-SSRS strict):} ``Critical to avoid false positives.''
Required \emph{persistent} ideation.
Result: F1=0.793, Recall=0.738. \textit{Too conservative.}

\textbf{v2 (err-on-caution):} Removed false-positive framing,
lowered the ideation bar to single expressions.
Result: F1=0.838, Recall=\textbf{0.877}. \textit{Best recall; low precision.}

\textbf{v3 (context-aware):} Added stressor-presence rule.
Ambiguous phrases safe only when a clear named stressor
explains them. \textit{Evaluation pending.}

\subsection{Key Design Decisions}

\textbf{Confidence override} at 0.85 bypasses the blend
for high-conviction LLM signals (upward only --- safety-first).

\textbf{Threshold 0.40} biases toward recall.
Consistent across all v3 runs including the 0.880 best run.

\textbf{Deterministic DDM} via \texttt{seed=42} (commit \texttt{d74c879})
eliminates Gaussian noise run-to-run variance.

\subsection{Known Code Gaps}\label{sec:gaps}

The following features exist in the codebase but are not connected:

\begin{itemize}\setlength\itemsep{0pt}
  \item \texttt{\_detect\_lang()} in \texttt{fr\_patch.py} --- never called
  \item \texttt{score\_with\_convokit()} --- implemented, not invoked
  \item \texttt{\_coordination\_score()} --- computed but not in WEIGHTS
  \item Output guardrail always \texttt{None} --- no response-side filter
  \item No conversation \texttt{history} passed in \texttt{context} ---
        DDM sees a string, not incremental turns
  \item LLM failure falls back silently to Cassandre-only scoring
\end{itemize}

%% ── 4. RED-TEAM FINDINGS ────────────────────────────────────────
\section{Red-Team / Stress-Test Findings}

\subsection{Methodology}

Two rounds of stress-testing against the live KHP chatbot:
Round v2 --- 92 handcrafted scenarios;
Round v2.1 --- 1,977 DEI-stratified LLM-generated scenarios
with automated judging (\texttt{auto\_judge\_v2.py}).

\subsection{Baseline Chatbot Failures}

\begin{center}
\captionof{table}{KHP critical miss rates (v2.1, $n$=1977)}\label{tab:failures}
\small
\begin{tabular}{lcc}
\toprule
\textbf{Category} & \textbf{Crit.\ miss} & \textbf{Pass}\\
\midrule
Hate speech (HR) & 97.7\% & 2.3\%\\
Body image (HR) & 83.7\% & 2.3\%\\
Home violence & 81.4\% & 9.3\%\\
Identity threats & 81.4\% & 16.3\%\\
Grief + SI & 67.4\% & 32.6\%\\
Suicide explicit & 0\% & \textbf{100\%}\\
\bottomrule
\end{tabular}
\end{center}

The chatbot handles explicit suicide mentions well but catastrophically fails
on indirect signals --- exactly where a guardrail adds value.

\subsection{DEI Robustness Gap}

Pass rates by persona:
child 10yr: 77.8\%;
neurodivergent: 73.1\%;
2SLGBTQ+: 73.7\%;
university student: 66.7\%;
\textbf{racialized minority: 58.8\%}.
The 19-point gap for racialized youth suggests systematic under-detection
of culturally-specific distress expression.

\subsection{Annotated Failures}

\begin{redbox}\small
\textbf{Miss --- Home Violence (racialized, 17F):}\\
\textit{User:} ``My dad hits me when he drinks.
It happens a lot. I just stay in my room.''\\
\textit{KHP:} ``That sounds stressful. Would you like to talk to a
counsellor about managing stress at home?''\\
\textit{Failure:} Reframed active abuse as stress.
No escalation. Our guardrail: $s=0.87 \to$ \texttt{HIGH\_RISK}.
\end{redbox}

\begin{amberbox}\small
\textbf{Miss --- Grief + SI (QC French, neurodivergent):}\\
\textit{User:} ``Je pense souvent \`a mon fr\`ere qui est mort.
Des fois je me dis que \c{c}a serait correct que je sois avec lui.''\\
\textit{KHP:} ``C'est normal de penser \`a ceux qu'on a perdus.''\\
\textit{Failure:} Missed passive SI (\textit{\^etre avec lui}).
Our FR patch catches the resignation tone; DDM detects flat affect.
\end{amberbox}

%% ── 5. QUANTITATIVE RESULTS ─────────────────────────────────────
\section{Quantitative Performance}

\begin{center}
\captionof{table}{Evaluation history, hidden-validation-102}\label{tab:results}
\small
\begin{tabular}{lcccc}
\toprule
\textbf{Run} & \textbf{F1} & \textbf{P} & \textbf{R} & \textbf{ms}\\
\midrule
Best (v2.5) & \textbf{.880} & \textbf{.917} & .846 & 1128\\
v3 few-shot & .850 & .871 & .831 & 1189\\
Bilingual v1 & .793 & .857 & .738 & 1074\\
Bilingual v2 & .838 & .803 & \textbf{.877} & 1289\\
\bottomrule
\end{tabular}
\end{center}

\textbf{Key finding:} There is a fundamental precision-recall trade-off in prompt design.
Conservative prompts maximise precision but miss genuine crises.
Aggressive prompts achieve highest recall but generate false alarms.
The v3 context-aware prompt attempts to resolve this through
stressor-presence reasoning rather than a blunt threshold.

\textbf{Latency:} 1,050--1,290\,ms/sample on the hidden set.
Acceptable for asynchronous triage; borderline for real-time turn-by-turn.

%% ── 6. KHP USABILITY ────────────────────────────────────────────
\section{KHP Usability \& De-escalation}

The guardrail operates as an \textbf{input filter}:
it evaluates the user message before the chatbot responds.
A \texttt{HIGH\_RISK} flag would trigger escalation to a human counsellor.

\textbf{Alert fatigue.} At precision 0.803, $\sim$1 in 5 flags is a false positive.
The context-aware v3 prompt is designed to address this.

\textbf{Latency.} 1.1--1.3\,s per turn is noticeable in a chat interface.
Acceptable for first-turn triage; may degrade mid-conversation UX.

\textbf{Bilingual coverage.} The system correctly handles QC French joual,
code-switching, and youth slang.
31\% of evaluation data is French or mixed-language --- a critical production requirement.

%% ── 7. DEPLOYMENT READINESS ─────────────────────────────────────
\section{Deployment Readiness}

\begin{center}\small
\begin{tabular}{cl}
{\color{safegreen}\textbf{\checkmark}} & Bilingual EN/FR with QC joual\\
{\color{safegreen}\textbf{\checkmark}} & Multi-provider LLM fallback\\
{\color{safegreen}\textbf{\checkmark}} & Deterministic DDM (seed=42)\\
{\color{safegreen}\textbf{\checkmark}} & ADHD/ASD profile detection\\
{\color{safegreen}\textbf{\checkmark}} & BUZZ-only runtime (no external APIs)\\
{\color{safegreen}\textbf{\checkmark}} & Submission contract satisfied\\
{\color{riskred}\textbf{\texttimes}} & No output guardrail\\
{\color{riskred}\textbf{\texttimes}} & No production monitoring\\
{\color{riskred}\textbf{\texttimes}} & Stateless (no session history)\\
{\color{riskred}\textbf{\texttimes}} & No rate limiting or abuse detection\\
{\color{riskred}\textbf{\texttimes}} & Silent LLM fallback\\
{\color{riskred}\textbf{\texttimes}} & No human escalation integration\\
\end{tabular}
\end{center}

%% ── 8. RISKS, LIMITATIONS, FUTURE WORK ─────────────────────────
\section{Risks, Limitations \& Future Directions}

\subsection{Technical Limitations}
The guardrail has no memory of prior sessions.
A user who disclosed suicidal ideation previously
receives no elevated prior in a new session.

The LLM judge is non-deterministic.
Despite a fixed DDM seed, two identical conversations
can receive different risk scores.

The 23-category taxonomy provides DEI coverage in generation,
but the evaluation set is dominated by two categories.
Performance on rare classes (Hate Speech: $n=1$; Grief: $n=1$)
cannot be reliably estimated.

\subsection{Future Improvements}
\textbf{Fine-tuned classifier.}
Replace or augment the LLM judge with a fine-tuned
mmBERT/CamemBERT classifier. Eliminates LLM latency and non-determinism.

\textbf{Stateful risk accumulation.}
Pass conversation history through \texttt{context}
across API calls for true session-level DDM accumulation.

\textbf{Linguistic coordination.}
Wire \texttt{score\_with\_convokit()} to detect user-bot
vocabulary mirroring as a dependency signal.

\textbf{Output guardrail.}
Implement response-side filtering to prevent harmful
de-escalation language (see home violence example).

\textbf{Threshold calibration.}
Optimise the composite threshold (0.40) using
precision-recall curves on a held-out set.

%% ── 9. CONCLUSION ───────────────────────────────────────────────
\section{Conclusion: On the Fundamental Limits of AI in Youth Mental Health}

This work demonstrates that a well-engineered guardrail
can meaningfully improve the safety of a youth mental health chatbot.
Our system achieves recall of 0.877 on a challenging bilingual evaluation set,
and our red-team analysis reveals systematic failure modes in the KHP baseline
that a guardrail can partially mitigate.

\textbf{We conclude, however, that AI chatbots are a fundamentally inadequate
--- and potentially harmful --- intervention for youth in mental health crisis,
regardless of guardrail sophistication.}
This is not a technical limitation we expect to engineer away.
It is a structural one.

\textbf{Evidence of harm.}
Miner et al.\ \cite{miner2016} demonstrated that major conversational agents
(Siri, Google Now, Cortana) failed to respond appropriately to suicidal disclosures,
sometimes responding with jokes or irrelevant information.
More recently, the National Eating Disorders Association's chatbot ``Tessa''
was withdrawn in 2023 after providing caloric restriction advice
to eating disorder patients \cite{neda2023} ---
the direct opposite of safe clinical guidance.
These are not edge cases of poor implementation.
They reflect the fundamental gap between pattern matching and clinical judgment.

\textbf{Therapeutic alliance cannot be replicated.}
Decades of psychotherapy research establish that the therapeutic alliance ---
the quality of the human relationship between clinician and client ---
is the strongest predictor of treatment outcomes,
exceeding the effect of any specific intervention \cite{horvath1991}.
Baumeister and Leary's foundational work \cite{baumeister1995}
establishes genuine human connection as a basic psychological need.
A chatbot cannot form an alliance.
It cannot be genuinely present.
It cannot follow up.
It cannot call for help.

\textbf{Youth are especially vulnerable.}
Adolescent identity formation and emotional regulation
are developmental processes that unfold in \emph{relationships}.
Orben and Przybylski \cite{orben2019} demonstrate that
technology use is associated with adolescent wellbeing changes.
Directing distressed youth toward AI systems
risks substituting a simulacrum of care for the real thing ---
one that is always available, never tired,
and entirely without the capacity to act when a child's life is at risk.

\textbf{A guardrail is not a safety net.}
The most important finding of our red-team analysis is not our recall score.
It is that the KHP chatbot missed 81--97\% of high-risk cases
in several categories.
A guardrail is downstream of a system that should not be
the primary point of contact for a child in crisis.
A child who discloses abuse and receives a response
about ``managing stress at home'' has not been helped.
They have been failed by a system that had the form of support
but not the substance.

\textbf{The appropriate use of this work}
is as a triage layer that routes high-risk signals
to human counsellors faster.
That is a meaningful, bounded, and honest contribution.
It is the only one technology can reliably make here.

\begin{redbox}
\small\textbf{Our position:}
No guardrail architecture, however robust,
justifies deploying an AI chatbot as a primary mental health intervention
for children. The right question is not
``how good is the guardrail?'' but ``why is there not a human here?''
\end{redbox}

\end{multicols}

%% ── REFERENCES ──────────────────────────────────────────────────
\begin{thebibliography}{99}\small\setlength\itemsep{1pt}
\bibitem{ratcliff2008} Ratcliff, R., \& McKoon, G. (2008). The diffusion decision model. \textit{Neural Computation}, 20(4), 873--922.
\bibitem{karalunas2012} Karalunas, S.~L., et al. (2012). Decomposing ADHD-related effects. \textit{Neuropsychology}, 26(6), 684.
\bibitem{south2010} South, M., et al. (2010). Cognitive flexibility in ASD. \textit{Journal of Neuroscience}.
\bibitem{miner2016} Miner, A.~S., et al. (2016). Smartphone-based conversational agents and responses to questions about mental health. \textit{JAMA Internal Medicine}, 176(5), 619--625.
\bibitem{neda2023} National Eating Disorders Association. (2023). Statement on the suspension of Tessa chatbot. \url{nedawareness.org}.
\bibitem{horvath1991} Horvath, A.~O., \& Symonds, B.~D. (1991). Relation between working alliance and outcome in psychotherapy. \textit{Journal of Counseling Psychology}, 38(2), 139.
\bibitem{baumeister1995} Baumeister, R.~F., \& Leary, M.~R. (1995). The need to belong. \textit{Psychological Bulletin}, 117(3), 497.
\bibitem{orben2019} Orben, A., \& Przybylski, A.~K. (2019). Adolescent well-being and digital technology use. \textit{Nature Human Behaviour}, 3(2), 173--182.
\bibitem{fitzpatrick2017} Fitzpatrick, K.~K., et al. (2017). Delivering CBT using a fully automated conversational agent (Woebot). \textit{JMIR Mental Health}, 4(2), e19.
\end{thebibliography}

\end{document}
